\begin{ledgroupsized}[r]{120mm}
\footnotesize
\pstart
\noindent\textbf{Überlieferung:}
\pend
\end{ledgroupsized}
\begin{ledgroupsized}[r]{114mm}
\footnotesize
\pstart \parindent -6mm
\makebox[6mm][l]{\textit{L}}%
Aufzeichnung:
LH~XXXVII~4 Bl.~168. 
Ein als Schreibblatt wiederverwendeter Teil eines Briefumschlags (ca.~20~x~13,5~cm.);
Papiererhaltungsmaßnahmen;
oberer und linker Rand beschnitten.
Eine Seite auf Bl.~168~r\textsuperscript{o};
Bl.~\href{\Scan{LH_37_04_168v}}{168~v\textsuperscript{o}}
überliefert zwei Zeichnungen ohne Zusammenhang mit dem Text, sonst leer.
Auf Bl.~168~r\textsuperscript{o} befinden sich Briefsiegelreste sowie, in gegenläufiger Schreibrichtung, die Anschrift 
(Hand von \indexN{\textls{Bodenhausen} (Bodenus, Bodenausen), Rudolf Christian, Freiherr v. 1698}R.\,C. von Bodenhausen):	
\textit{All'\;Ill:\textsuperscript{mo} Sig.\textsuperscript{re} mio Sig.\textsuperscript{re} P\textsuperscript{\,rone} Col.\textsuperscript{mo}} \lbrack/\rbrack\
\textit{Il Sig.\textsuperscript{re} Gotofredo Gul.\textsuperscript{mo} Leibnitz, Consigl.\textsuperscript{ro}} \lbrack/\rbrack\
\textit{di S.A.S. di Bransvic.} \lbrack/\rbrack\ \textit{a} \lbrack/\rbrack\ \textit{Hannover}.
Der Träger bildete ursprünglich eine Einheit mit 
LBr~79 Bl.~31 (\cite{02131}Brief von \indexN{\textls{Bodenhausen} (Bodenus, Bodenausen), Rudolf Christian, Freiherr v. 1698}R.\,C. Bodenhausen an Leibniz vom 23. September 1690,
\textit{LSB} III, 4 N.\;275).	
\pend
\end{ledgroupsized}
\vspace{5mm}
\begin{ledgroup}
\footnotesize
\pstart
\noindent%
\textbf{Datierungsgründe:}
Leibniz hat die Aufzeichnung N.\;\ref{53450} auf der recto-Seite 
eines ursprünglich als Briefumschlag verwendeten, 
bis auf die Anschrift in R.\,C. Bodenhausens Hand nicht beschriebenen Blattes italienischen Papiers
verfasst.
Der dazugehörige Brief \indexN{\textls{Bodenhausen} (Bodenus, Bodenausen), Rudolf Christian, Freiherr v. 1698}Bodenhausens lässt sich auf Grund der gleichen Maße, Papiereigenschaften und der gemeinsamen Trennkante mit dem Träger LBr~79 Bl.~31 mit Sicherheit als
sein \cite{02131}Schreiben aus Florenz vom 23. September 1690 
(\textit{LSB} III, 4 N.\;275) identifizieren.
\pend
\pstart
Die zwei auf Bl.~168~v\textsuperscript{o} überlieferten,
keinen erkennbaren Zusammenhang mit N.\;\ref{53450} aufweisenden
Zeichnungen (Teildarstellungen einer Hyperbel und einer Zykloide)
sind ebenfalls mit dem genannten Brief in Beziehung zu setzen:
Leibniz beantwortet dieses und drei andere %
vorangegangene Schreiben Bodenhausens 
mit einem \cite{02132}Brief vom 26.~Oktober (5.~November) 1690 (\textit{LSB} III, 4 N.\;285),
dem er eine \cite{02133}\textit{scheda} mit Beispielen für die Leistungsfähigkeit des Calculus beilegt
(ebd. N.\;286).
Beide Diagramme finden sich darin in ausgearbeiteter Form wieder und dienen dort zur Veranschaulichung der Leibniz'schen Rektifikation der Hyperbel (S.\;638) bzw. der Zykloide (S.\;640).
Aus diesem Grund werden die Diagramme nicht mit N.\;\ref{53450} ediert.
Die Entstehung der Zeichnungen lässt sich daher auf den Zeitraum Ende September bis 5. November einkreisen.
\pend
\pstart
Der Zeitpunkt des Erhalts des Briefs durch Leibniz (wahrscheinlich Ende September 1690) bietet einen Terminus post quem für die Entstehung von N.\;\ref{53450}. 
Unter der Annahme einer zeitnahen Beschreibung des Blattes, die durch die sichere Entstehung der Diagramme auf Bl.~168~v\textsuperscript{o} vor dem 5. November 1690 bekräftigt wird,
ergibt sich die vorgeschlagene Datierung.
\pend
\end{ledgroup}
\newpage
\selectlanguage{latin}
\vspace{8mm}
\pstart%
\noindent%
\href{\Scan{LH_37_04_168r}}%
{\lbrack168~r\textsuperscript{o}\rbrack}
\pend
\vspace{1.0em} %
\centerline{\includegraphics[width=0.165\textwidth]{%
gesamttex/edit_VIII-4/images/LH_37_04_168_d1_168r.pdf}} 
\nobreak\vspace{0.5em}
\centerline{\lbrack\textit{Fig.~1}\rbrack}
\vspace{1.0em}
\pstart
\edtext{}{%
\lemma{\lbrack\textit{Fig.~1}\rbrack}%
\Cfootnote[nonum]{%
Ein gestrichener und von Leibniz überschriebener Entwurf zum Diagramm wird nicht wiedergegeben.}}%
Si corpora \textit{A} et \textit{B}
\edtext{concurrant,\indexS{concursus} ita ut}{\lemma{concurrant,}\Bfootnote{%
\textit{(1)}~ut %
\textit{(2)}~ita ut~\textit{L}}}
\textit{B} et \textit{C} intelligantur eandem semper servare distantiam\indexS{distantia eadem servata} inter se, et eundem semper 
angulum \textit{ABC}; id est ut vel quiescant simul \textit{B} et \textit{C} vel rectas parallelas describant, 
oportet ut discerni non possit per eventum
\edtext{concursus\indexS{concursus}\indexS{eventus concursus} quantus}{\lemma{concursus}\Bfootnote{%
\textit{(1)}~quantus cuique sit %
\textit{(2)}~utrum motus %
\textit{(3)}~quantus~\textit{L}}}
motus\indexS{motus} sit in altero, aut in quo sit quies\indexS{quies} inter \textit{B} et \textit{C}. 
Sit centrum\indexS{centrum gravitatis} \textit{G}, debet celeritas respectiva\indexS{celeritas respectiva}\indexS{celeritas respectiva eadem manens} eadem manere inter \textit{A}
\edtext{et \textit{G}. Quod}{\lemma{et \textit{G}.}\Bfootnote{%
\textit{(1)}~Sit porro centrum inter %
\textit{(2)}~Quod~\textit{L}}}
fiet,
\edtext{sive liberum}{\lemma{sive}\Bfootnote{%
\textit{(1)}~sint %
\textit{(2)}~liberum~\textit{L}}}
sit \textit{B} a \textit{C}, sive sint
connexa.\indexS{corpora connexa} 
Jam si liberum esset \textit{B} a \textit{C} orietur vis respectiva\indexS{vis respectiva} inter \textit{B}
\edtext{et \textit{C}. Hanc}{\lemma{et \textit{C}}\Bfootnote{%
\textit{(1)}~ea scilicet quae %
\textit{(2)}~. Hanc~\textit{L}}}
eandem vim respectivam servari\indexS{vis respectiva}\indexS{conservatio vis respectivae} 
necesse est inter \textit{B} et \textit{C} etsi sint
\edtext{connexa per circulationem\indexS{circulatio}\indexS{circulatio circa centrum} circa}{\lemma{connexa}\Bfootnote{%
\textit{(1)}~. Adeoque %
\textit{(2)}~. Unde %
\textit{(3)}~per circulationem circa~\textit{L}}}
centrum\indexS{centrum gravitatis} \textit{G}. 
Inter \textit{A} et \textit{B} sit centrum\indexS{centrum gravitatis} \textit{H}, et inter \textit{A} et \textit{C} sit centrum \textit{L}. 
Oportet ut semper eadem sit vis respectiva\indexS{vis respectiva}\indexS{vis respectiva eadem manens} inter 
\edtext{\lbrack\textit{A}\rbrack}{%
\lemma{}%
\Bfootnote{%
\textit{H} %
\textit{L ändert Hrsg.}}}
et \textit{G}, inter \textit{B} et \textit{L}\lbrack,\rbrack\ inter \textit{C} et \textit{H}. 
\pend
\pstart
Videndum an hoc fiat et in circulatione\indexS{circulatio} dicto modo determinata et an ex hoc solo de servata eadem vi respectiva\indexS{vis respectiva}\indexS{vis respectiva eadem servata} determinetur
\edtext{circulatio,\indexS{circulatio} posito servari eandem vim respectivam\indexS{vis respectiva}\indexS{vis respectiva eadem servata} et eandem inter \textit{B} et \textit{C} distantiam.\indexS{distantia}\indexS{distantia eadem servata}}{\lemma{circulatio}\Bfootnote{%
\textit{(1)}~. Sed quaeritur de ipsa celeritate circulationis. Sane illud discerni potest ex phaenomenis utrum %
\textit{(2)}~, posito servari \lbrack...\rbrack\ distantiam.~\textit{L}}}
Videndum an per circulationem\indexS{circulatio} 
translato \textit{B} in (\textit{B}) et \textit{C}
\edtext{in (\textit{C}) ipso}{%
\lemma{}%
\Bfootnote{%
in (\textit{C}) \bars{adhuc \textit{gestr.}} %
ipso~\textit{L}}}
autem \textit{A} in recta procedente, servetur 
vis respectiva\indexS{vis respectiva}\indexS{vis respectiva eadem servata} inter \textit{B} et \textit{L}, 
seu inter \textit{C} et \textit{H} quia servatur inter 
\edtext{\lbrack\textit{A}\rbrack}{%
\lemma{}%
\Bfootnote{%
\textit{H} %
\textit{L ändert Hrsg.}}}
et \textit{G}.
\pend
\pstart 
Examinandum etiam quid
\edtext{dicendum}{\lemma{}\Bfootnote{dicendum \textit{erg.~L}}}
\edtext{de circulatione,\indexS{circulatio}}{%
\lemma{de}%
\Bfootnote{%
\textit{(1)}~circulationibus, %
\textit{(2)}~circulatione,~\textit{L}}}
quatenus ex his phaenomenis cognosci possit centrum\indexS{centrum circulationis} ejus, 
et vertiginis\indexS{vertigo} quantitas.\indexS{quantitas vertiginis}
Existenti 
\edtext{in \textit{A} examinandum quid}{\lemma{in \textit{A}}\Bfootnote{%
\textit{(1)}~\textit{G} semper movebitur in linea recta, sed vi %
\textit{(2)}~examinandum quid~\textit{L}}}
appareat et quid ex ipsius apparentia\indexS{apparentia} ante et post concursum\indexS{concursus}
possit determinari, 
similiter quid ex \textit{B}, vel ex \textit{C} et quid prodeat,
\edtext{si habitans}{\lemma{si}\Bfootnote{%
\textit{(1)}~unusquisque se ponat %
\textit{(2)}~habitans~\textit{L}}}
cujusque corporis se ponat quiescentem. 
\pend
\pstart 
In motu\indexS{motus} examinemus, quae constantia sine concursu,\indexS{concursus}
quibus conjungenda
\edtext{deinde quidem manet quoque eadem quantitas}{\lemma{deinde}\Bfootnote{%
\bars{quidem \textit{erg.}} manet %
\textit{(1)}~eadem quantitas quoque %
\textit{(2)}~quoque eadem quantitas~\textit{L}}}
motus\indexS{quantitas motus}\indexS{quantitas motus eadem manens}
cum nullus est concursus.\indexS{concursus} Sed tamen hinc non sequitur
\edtext{manere et in concursu.\indexS{concursus}}{%
\lemma{manere}%
\Bfootnote{%
\textit{(1)}~\bars{et \textit{streicht Hrsg.}} donec %
\textit{(2)}~et in concursu.~\textit{L}}}
\pend
\pstart 
Optima Methodus
\edtext{ut ea mane\lbrack a\rbrack nt quae}{%
\lemma{ut}%
\Bfootnote{%
\textit{(1)}~ex %
\textit{(2)}~ea %
\bars{maneat \textit{ändert Hrsg.}} %
\textit{(a)}~etiam ante concursum %
\textit{(b)}~quae~\textit{L}}}
impenetrabilitas\indexS{impenetrabilitas}
\edtext{non}{\lemma{}\Bfootnote{non \textit{erg.~L}}}
impedit\lbrack,\rbrack\ si in astris\indexS{aster} concipiamus vim trajectoriam\indexS{vis trajectoria}
et simul vim magneticam\indexS{vis magnetica} 
seu respectivam\indexS{vis respectiva}\lbrack,\rbrack\ vel potest definiri in
\edtext{quo quies.\indexS{quies}}{\lemma{quo}\Bfootnote{%
\textit{(1)}~motus %
\textit{(2)}~quies.~\textit{L}}}
\pend
\pstart
Sed ratio\indexS{ratio pro quiete solis} pro quiete\indexS{quies solis} hac solis\indexS{sol} 
haec est, quod motus 
qui isti tribuitur rectilineus\indexS{motus rectilineus}\indexS{motus solis rectilineus} tribuendus esset 
omnibus fixis.\indexS{stellae fixae}\indexS{sidera fixa}
\pend
\pstart 
Omnis motus rectilineus.\indexS{motus rectilineus} 
Curvilinei\indexS{motus curvilineus} non nisi compositi.\indexS{motus compositus}
Hinc omnis firmitas\indexS{firmitas} non est nisi motus unius ad alterum 
ab ambiente\indexS{ambiens} impressus.\indexS{motus ab ambiente impressus}
\pend 
